\documentclass[DIN, pagenumber=false, fontsize=11pt, parskip=half]{scrartcl}
\usepackage[UTF8]{ctex}  % 中文支持
% \usepackage{ngerman}
\usepackage[utf8]{inputenc}
\usepackage[T1]{fontenc}
\usepackage{textcomp}
\usepackage{graphicx}
\usepackage{amsmath}
\usepackage{booktabs}
\usepackage{tasks}
\usepackage{xcolor} % Required for custom red

\usepackage{geometry} % Custom margins

\usepackage{tocloft} % Modifying the table of contents

% for matlab code
% bw = blackwhite - optimized for print, otherwise source is colored
\usepackage[framed,numbered,bw]{mcode}
\usepackage[most]{tcolorbox}
% for other code
%\usepackage{listings}

\setlength{\parindent}{0em}

% set section in CM
\setkomafont{section}{\normalfont\bfseries\Large}

\newcommand{\mytitle}[1]{{\noindent\Large\textbf{#1}}}
\newcommand{\mysection}[1]{\textbf{\section*{#1}}}
\newcommand{\mysubsection}[2]{\romannumeral #1) #2}

\newtcolorbox{parchmentquote}{
    enhanced,
    colback=orange!5!white,   % 浅米黄色背景
    colframe=orange!50!black,
    boxrule=0.5pt,
    arc=2pt,                  % 细微圆角
    boxsep=5pt,
    before skip=10pt,
    after skip=10pt,
    fontupper=\itshape,
    borderline west={2pt}{0pt}{orange!80!black} % 加厚的左装饰线
}
%===================================
\begin{document}

\noindent\textbf{课后练习} \hfill \textbf{online course}\\
primary school \hfill Liu\\

\mytitle{比例尺 \hfill \today}
\newline
\begin{parchmentquote}
    本节课主要复习比例尺的相关知识。
    从基础开始，学习比例尺的定义、应用以及相关的计算问题。


\end{parchmentquote}

\section{什么是比例尺？}

生活中，真实物体太大（或太小），不能直接画出来，
所以需要按照相同的倍数缩小（或放大）。

例如：

\begin{itemize}
    \item 地图
    \item 小区平面图
    \item 房屋设计图
    \item 教室座位图
\end{itemize}

\textbf{比例尺就是：}

\[
\text{图上距离} : \text{实际距离}
\]

例如：

图上 $1$ 厘米表示实际 $100$ 厘米，

比例尺写成：

\[
1:100
\]

读作：

\textbf{比例尺一比一百。}

\section{为什么需要比例尺？}

例如：

黑板长 $4$ 米，

如果直接画，纸放不下。

规定：

\[
1\text{ cm 表示 }1\text{ m}
\]

那么：

\[
4\text{ m}\rightarrow 4\text{ cm}
\]

这就是按同样倍数缩小。

\section{比例尺的计算}

\subsection{求实际距离}

已知：

\begin{itemize}
    \item 比例尺：$1:50000$
    \item 图上距离：$3$ cm
\end{itemize}

计算：

\[
3\times 50000=150000\text{ cm}
\]

换算：

\[
150000\text{ cm}
=
1500\text{ m}
=
1.5\text{ km}
\]

\textbf{口诀：}

\[
\boxed{\text{求实际：乘分母}}
\]

\subsection{求图上距离}

已知：

\begin{itemize}
    \item 实际距离：$800$ m
    \item 比例尺：$1:40000$
\end{itemize}

先统一单位：

\[
800\text{ m}=80000\text{ cm}
\]

计算：

\[
80000\div40000=2\text{ cm}
\]

\textbf{口诀：}

\[
\boxed{\text{求图上：实际除分母}}
\]

\subsection{求比例尺}

已知：

\begin{itemize}
    \item 图上距离：$5$ cm
    \item 实际距离：$100$ m
\end{itemize}

先统一单位：

\[
100\text{ m}=10000\text{ cm}
\]

写成比：

\[
5:10000
\]

化简：

\[
1:2000
\]

\textbf{口诀：}

\[
\boxed{\text{求比例尺：先同单位，再化简}}
\]

\section{练习}

\begin{enumerate}
    \item 比例尺为 $1:1000$，图上距离为 $5$ cm，
    求实际距离。

    \item 实际距离为 $300$ m，比例尺为 $1:10000$，
    求图上距离。

    \item 图上距离为 $4$ cm，实际距离为 $80$ m，
    求比例尺。
\end{enumerate}

\section{容易出错的地方}

\begin{enumerate}
    \item 单位必须统一。

    例如：

    \[
    2\text{ m}=200\text{ cm}
    \]

    所以比例尺是：

    \[
    1:200
    \]

    不能直接写成：

    \[
    1:2
    \]

    \item 比例尺不能带单位。

    错误：

    \[
    1\text{ cm}:100\text{ cm}
    \]

    正确：

    \[
    1:100
    \]

    \item 比例尺前项通常化成 $1$。
\end{enumerate}

\section{课堂活动}

请画出自己的房间平面图。

规定：

\[
1\text{ cm 表示 }50\text{ cm}
\]

测量：

\begin{itemize}
    \item 床的长度
    \item 书桌长度
    \item 门和窗的位置
\end{itemize}

最后画出房间平面图。

体会：

\[
\boxed{\text{比例尺就是按相同倍数缩小}}
\]

\end{document}